\documentclass[11pt]{article}
\usepackage[T1]{fontenc}
\usepackage{textcomp}
\usepackage[margin=0.75in]{geometry}
\usepackage{amsmath,amssymb,amsthm}
\usepackage{array,booktabs}
\usepackage{hyperref}
\setlength{\parindent}{0pt}
\newtheorem{theorem}{Theorem}
\theoremstyle{definition}
\newtheorem{definition}{Definition}
\title{Flow Proof Artifact: prop-05-rectangle-unequal-with-square-between}
\author{Generated by \texttt{flow doc proof}}
\date{Flow Proof Book}
\begin{document}
\maketitle
If a straight line is cut into equal and unequal segments, the rectangle by the unequal segments together with the square on the line between the points of section equals the square on the half.\\[0.5em]
\textbf{Source.} euclid — Elements, Book II, Proposition 5\\[0.5em]
\bigskip
\noindent{\large \textbf{Derived fact 1.} \textit{If a straight line is cut into equal and unequal segments, the rectangle by the unequal segments together with the square on the line between the points of section equals the square on the half} \hfill the Euclidean plane \cdot Euclid Book II \cdot Proposition 5: the rectangle by unequal parts plus the square on the midpoint equals the square on the half}\par
\smallskip\noindent\textit{Coordinate: the Euclidean plane \cdot Euclid Book II \cdot Proposition 5: the rectangle by unequal parts plus the square on the midpoint equals the square on the half}\par
\smallskip\noindent\textit{Source: euclid — Elements, Book II, Proposition 5}\par
\smallskip\noindent\textit{Built on: proposition 4: the square on the whole equals the squares on the parts plus twice their rectangle, for Euclid Book II on the Euclidean plane}\par
\medskip
\noindent\textbf{Goal.} If a straight line is cut into equal and unequal segments, the rectangle by the unequal segments together with the square on the line between the points of section equals the square on the half\par
\noindent$\displaystyle rect(\text{unequal parts}) + sq(\text{between points}) = sq(half)$\par
\medskip
\noindent
\begin{tabular}{@{} >{\raggedright\arraybackslash}p{0.06\textwidth} >{\raggedright\arraybackslash}p{0.47\textwidth} >{\raggedleft\arraybackslash}p{0.41\textwidth} @{}}
\textbf{\#} & \textbf{Proof} & \textbf{Mathematics} \\
\midrule
\textbf{1.} & We prove that proposition 5: the rectangle by unequal parts plus the square on the midpoint equals the square on the half for Euclid Book II the Euclidean plane. & \\[0.45em]
\textbf{2.} & Let whole = straight\_line\_AB. & \\[0.45em]
\textbf{3.} & Let half = segment\_AM. & \\[0.45em]
\textbf{4.} & Let unequal = segments\_AN\_and\_NB. & \\[0.45em]
\textbf{5.} & From step 2, step 3, and step 4, we can deduce that rect(unequal) plus sq(between points) equals sq(half). & $\displaystyle rect(unequal) + sq(\text{between points}) = sq(half)$ \\[0.45em]
\textbf{6.} & We invoke the derived fact governing Euclid Book II the Euclidean plane: proposition 4: the square on the whole equals the squares on the parts plus twice their rectangle, for Euclid Book II on the Euclidean plane. & \\[0.45em]
\textbf{7.} & From step 6, this implies rect(unequal parts) plus sq(between points) equals sq(half). Hence proven. & $\displaystyle rect(\text{unequal parts}) + sq(\text{between points}) = sq(half)$ \\[0.45em]
\bottomrule
\end{tabular}
\label{thm:Geometry-Euclid Book II-proposition_5_the_rectangle_by_unequal_parts_plus_the_square_on_the_midpoint_equals_the_square_on_the_half}
\par\medskip\hrule
\end{document}
